\documentclass{amsart}
% For dealing with references we use the comment environment
\usepackage{verbatim}
\newenvironment{reference}{\comment}{\endcomment}
%\newenvironment{reference}{}{}
\newenvironment{slogan}{\comment}{\endcomment}
\newenvironment{history}{\comment}{\endcomment}

% For commutative diagrams we use Xy-pic
\usepackage[all]{xy}

% We use 2cell for 2-commutative diagrams.
\xyoption{2cell}
\UseAllTwocells

% We use multicol for the list of chapters between chapters
\usepackage{multicol}

% This is generally recommended for better output
\usepackage{lmodern}
\usepackage[T1]{fontenc}

% For cross-file-references

% Package for hypertext links:
\usepackage{hyperref}

% For any local file, say "hello.tex" you want to link to please

% Theorem environments.
%
\theoremstyle{plain}
\newtheorem{theorem}[subsection]{Theorem}
\newtheorem{proposition}[subsection]{Proposition}
\newtheorem{lemma}[subsection]{Lemma}

\theoremstyle{definition}
\newtheorem{definition}[subsection]{Definition}
\newtheorem{example}[subsection]{Example}
\newtheorem{exercise}[subsection]{Exercise}
\newtheorem{situation}[subsection]{Situation}

\theoremstyle{remark}
\newtheorem{remark}[subsection]{Remark}
\newtheorem{remarks}[subsection]{Remarks}

\numberwithin{equation}{subsection}

% Macros
%
\def\lim{\mathop{\mathrm{lim}}\nolimits}
\def\colim{\mathop{\mathrm{colim}}\nolimits}
\def\Spec{\mathop{\mathrm{Spec}}}
\def\Hom{\mathop{\mathrm{Hom}}\nolimits}
\def\Ext{\mathop{\mathrm{Ext}}\nolimits}
\def\SheafHom{\mathop{\mathcal{H}\!\mathit{om}}\nolimits}
\def\SheafExt{\mathop{\mathcal{E}\!\mathit{xt}}\nolimits}
\def\Sch{\mathit{Sch}}
\def\Mor{\mathop{\mathrm{Mor}}\nolimits}
\def\Ob{\mathop{\mathrm{Ob}}\nolimits}
\def\Sh{\mathop{\mathit{Sh}}\nolimits}
\def\NL{\mathop{N\!L}\nolimits}
\def\CH{\mathop{\mathrm{CH}}\nolimits}
\def\proetale{{pro\text{-}\acute{e}tale}}
\def\etale{{\acute{e}tale}}
\def\QCoh{\mathit{QCoh}}
\def\Ker{\mathop{\mathrm{Ker}}}
\def\Im{\mathop{\mathrm{Im}}}
\def\Coker{\mathop{\mathrm{Coker}}}
\def\Coim{\mathop{\mathrm{Coim}}}

% Boxtimes
%
\DeclareMathSymbol{\boxtimes}{\mathbin}{AMSa}{"02}

%
% Macros for moduli stacks/spaces
%
\def\QCohstack{\mathcal{QC}\!\mathit{oh}}
\def\Cohstack{\mathcal{C}\!\mathit{oh}}
\def\Spacesstack{\mathcal{S}\!\mathit{paces}}
\def\Quotfunctor{\mathrm{Quot}}
\def\Hilbfunctor{\mathrm{Hilb}}
\def\Curvesstack{\mathcal{C}\!\mathit{urves}}
\def\Polarizedstack{\mathcal{P}\!\mathit{olarized}}
\def\Complexesstack{\mathcal{C}\!\mathit{omplexes}}
% \Pic is the operator that assigns to X its picard group, usage \Pic(X)
% \Picardstack_{X/B} denotes the Picard stack of X over B
% \Picardfunctor_{X/B} denotes the Picard functor of X over B
\def\Pic{\mathop{\mathrm{Pic}}\nolimits}
\def\Picardstack{\mathcal{P}\!\mathit{ic}}
\def\Picardfunctor{\mathrm{Pic}}
\def\Deformationcategory{\mathcal{D}\!\mathit{ef}}

\setlength{\emergencystretch}{3em}
\begin{document}
\title[Stacks Project, Tag 0AA0]{258 --- Compactification independence of the twisted inverse image functor}
\maketitle
\noindent Copyright (C) 2005--2025 Johan de Jong. Stacks Project authors. Source: \href{https://stacks.math.columbia.edu/tag/0AA0}{Tag 0AA0}.

\noindent Editorial difficulty note (not author text). Difficulty H2: The proof compares compactifications through dense common refinements and constructs comparison maps using proper-adjoint localization. It then proves independence of the common refinement and a cocycle identity for successive refinements. The complete construction of a canonical functor is the nonroutine obligation..

\noindent Editorial provenance note (not author text). History: excerpt from revision \texttt{a0444\allowbreak 6e57e\allowbreak c1fbc\allowbreak 252a8\allowbreak 71afc\allowbreak ec775\allowbreak 2fb28\allowbreak 07b14}, duality.tex, lines 3768--3941. Reference presentation was adapted; original-author context and core support blocks were retained or added. The mathematical wording is unchanged. Licensed under GNU FDL 1.2 or later; no invariant sections or cover texts. See ../licenses/COPYING. Context blocks reproduce author definitions and supporting results; they are not additional counted problems.

\begin{situation}
\label{situation-shriek}
Here $S$ is a Noetherian scheme and $\textit{FTS}_S$ is the category whose
\begin{enumerate}
\item objects are schemes $X$ over $S$ such that the structure
morphism $X \to S$ is both separated and of finite type, and
\item morphisms $f : X \to Y$ between objects are morphisms
of schemes over $S$.
\end{enumerate}
\end{situation}

\begin{lemma}
\label{lemma-shriek-well-defined}
In Situation \ref{situation-shriek} let $f : X \to Y$ be a morphism
of $\textit{FTS}_S$. The functor $f^!$ is, up to canonical isomorphism,
independent of the choice of the compactification.
\end{lemma}

\begin{proof}
The category of compactifications of $X$ over $Y$ is defined
in More on Flatness, Section \href{https://stacks.math.columbia.edu/tag/0ATT}{section compactify (Tag 0ATT)}.
By More on Flatness, Theorem \href{https://stacks.math.columbia.edu/tag/0F41}{theorem nagata (Tag 0F41)} and
Lemma \href{https://stacks.math.columbia.edu/tag/0A9Z}{lemma compactifyable (Tag 0A9Z)} it is nonempty.
To every choice of a compactification
$$
j : X \to \overline{X},\quad \overline{f} : \overline{X} \to Y
$$
the construction above associates the functor $j^* \circ \overline{a} :
D_\QCoh^+(\mathcal{O}_Y) \to D_\QCoh^+(\mathcal{O}_X)$
where $\overline{a}$ is the right adjoint of $R\overline{f}_*$
constructed in Lemma \href{https://stacks.math.columbia.edu/tag/0A9E}{lemma twisted inverse image (Tag 0A9E)}.

\medskip\noindent
Suppose given a morphism $g : \overline{X}_1 \to \overline{X}_2$
between compactifications $j_i : X \to \overline{X}_i$ over $Y$
such that $g^{-1}(j_2(X)) = j_1(X)$\footnote{This may fail
with our definition of compactification. See
More on Flatness, Section \href{https://stacks.math.columbia.edu/tag/0ATT}{section compactify (Tag 0ATT)}.}.
Let $\overline{c}$ be the right adjoint of
Lemma \href{https://stacks.math.columbia.edu/tag/0A9E}{lemma twisted inverse image (Tag 0A9E)} for $g$.
Then $\overline{c} \circ \overline{a}_2 = \overline{a}_1$
because these functors are adjoint to
$R\overline{f}_{2, *} \circ Rg_* = R(\overline{f}_2 \circ g)_*$.
By (\href{https://stacks.math.columbia.edu/tag/0A9L}{equation sheafy (Tag 0A9L)}) we have a canonical transformation
$$
j_1^* \circ \overline{c} \longrightarrow j_2^*
$$
of functors
$D^+_\QCoh(\mathcal{O}_{\overline{X}_2}) \to D^+_\QCoh(\mathcal{O}_X)$
which is an isomorphism by Lemma \href{https://stacks.math.columbia.edu/tag/0A9P}{lemma proper noetherian (Tag 0A9P)}.
The composition
$$
j_1^* \circ \overline{a}_1 \longrightarrow
j_1^* \circ \overline{c} \circ \overline{a}_2 \longrightarrow
j_2^* \circ \overline{a}_2
$$
is an isomorphism of functors which we will denote by $\alpha_g$.

\medskip\noindent
Consider two compactifications $j_i : X \to \overline{X}_i$, $i = 1, 2$
of $X$ over $Y$. By More on Flatness, Lemma
\href{https://stacks.math.columbia.edu/tag/0ATU}{lemma compactifications cofiltered (Tag 0ATU)} part (b)
we can find a compactification $j : X \to \overline{X}$
with dense image and morphisms
$g_i : \overline{X} \to \overline{X}_i$ of compactifications.
By More on Flatness, Lemma
\href{https://stacks.math.columbia.edu/tag/0ATU}{lemma compactifications cofiltered (Tag 0ATU)} part (c)
we have $g_i^{-1}(j_i(X)) = j(X)$. Hence we get isomorpisms
$$
\alpha_{g_i} :
j^* \circ \overline{a}
\longrightarrow
j_i^* \circ \overline{a}_i
$$
by the previous paragraph. We obtain an isomorphism
$$
\alpha_{g_2} \circ \alpha_{g_1}^{-1} :
j_1^* \circ \overline{a}_1 \to j_2^* \circ \overline{a}_2
$$
To finish the proof we have to show that these isomorphisms are well defined.
We claim it suffices to show the composition of isomorphisms constructed
in the previous paragraph is another (for a precise statement see the next
paragraph). We suggest the reader check this is true on a napkin, but we
will also completely spell it out in the rest of this paragraph.
Namely, consider a second choice of a compactification
$j' : X \to \overline{X}'$ with dense image
and morphisms of compactifications $g'_i : \overline{X}' \to \overline{X}_i$.
By More on Flatness, Lemma \href{https://stacks.math.columbia.edu/tag/0ATU}{lemma compactifications cofiltered (Tag 0ATU)}
we can find a compactification $j'' : X \to \overline{X}''$
with dense image and morphisms of compactifications
$h : \overline{X}'' \to \overline{X}$ and
$h' : \overline{X}'' \to \overline{X}'$. We may even assume
$g_1 \circ h = g'_1 \circ h'$ and $g_2 \circ h = g'_2 \circ h'$.
The result of the next paragraph gives
$$
\alpha_{g_i} \circ \alpha_h = \alpha_{g_i \circ h} =
\alpha_{g'_i \circ h'} = \alpha_{g'_i} \circ \alpha_{h'}
$$
for $i = 1, 2$. Since these are all isomorphisms of functors
we conclude that $\alpha_{g_2} \circ \alpha_{g_1}^{-1} =
\alpha_{g'_2} \circ \alpha_{g'_1}^{-1}$ as desired.

\medskip\noindent
Suppose given compactifications $j_i : X \to \overline{X}_i$
for $i = 1, 2, 3$.  Suppose given morphisms
$g : \overline{X}_1 \to \overline{X}_2$ and
$h : \overline{X}_2 \to \overline{X}_3$ of compactifications
such that $g^{-1}(j_2(X)) = j_1(X)$ and $h^{-1}(j_2(X)) = j_3(X)$.
Let $\overline{a}_i$ be as above. The claim above means that
$$
\alpha_g \circ \alpha_h = \alpha_{g \circ h} :
j_1^* \circ \overline{a}_1 \to j_3^* \circ \overline{a}_3
$$
Let $\overline{c}$, resp.\ $\overline{d}$ be the right adjoint of
Lemma \href{https://stacks.math.columbia.edu/tag/0A9E}{lemma twisted inverse image (Tag 0A9E)} for $g$, resp.\ $h$.
Then $\overline{c} \circ \overline{a}_2 = \overline{a}_1$ and
$\overline{d} \circ \overline{a}_3 = \overline{a}_2$
and there are canonical transformations
$$
j_1^* \circ \overline{c} \longrightarrow j_2^*
\quad\text{and}\quad
j_2^* \circ \overline{d} \longrightarrow j_3^*
$$
of functors
$D^+_\QCoh(\mathcal{O}_{\overline{X}_2}) \to D^+_\QCoh(\mathcal{O}_X)$
and
$D^+_\QCoh(\mathcal{O}_{\overline{X}_3}) \to D^+_\QCoh(\mathcal{O}_X)$
for the same reasons as above. Denote $\overline{e}$ the
right adjoint of Lemma \href{https://stacks.math.columbia.edu/tag/0A9E}{lemma twisted inverse image (Tag 0A9E)}
for $h \circ g$. There is a canonical transformation
$$
j_1^* \circ \overline{e} \longrightarrow j_3^*
$$
of functors
$D^+_\QCoh(\mathcal{O}_{\overline{X}_3}) \to D^+_\QCoh(\mathcal{O}_X)$
given by (\href{https://stacks.math.columbia.edu/tag/0A9L}{equation sheafy (Tag 0A9L)}). Spelling things out we have to
show that the composition
$$
\alpha_h \circ \alpha_g :
j_1^* \circ \overline{a}_1 \to
j_1^* \circ \overline{c} \circ \overline{a}_2 \to
j_2^* \circ \overline{a}_2 \to
j_2^* \circ \overline{d} \circ \overline{a}_3 \to
j_3^* \circ \overline{a}_3
$$
is the same as the composition
$$
\alpha_{h \circ g} :
j_1^* \circ \overline{a}_1 \to
j_1^* \circ \overline{e} \circ \overline{a}_3 \to
j_3^* \circ \overline{a}_3
$$
We split this into two parts. The first is to show that the diagram
$$
\xymatrix{
\overline{a}_1 \ar[r] \ar[d] & \overline{c} \circ \overline{a}_2 \ar[d] \\
\overline{e} \circ \overline{a}_3 \ar[r] &
\overline{c} \circ \overline{d} \circ \overline{a}_3
}
$$
commutes where the lower horizontal arrow comes from the identification
$\overline{e} = \overline{c} \circ \overline{d}$. This is true
because the corresponding diagram of total direct image functors
$$
\xymatrix{
R\overline{f}_{1, *} \ar[r] \ar[d] & Rg_* \circ R\overline{f}_{2, *} \ar[d] \\
R(h \circ g)_* \circ R\overline{f}_{3, *} \ar[r] &
Rg_* \circ Rh_* \circ R\overline{f}_{3, *}
}
$$
is commutative (insert future reference here). The second part
is to show that the composition
$$
j_1^* \circ \overline{c} \circ \overline{d} \to
j_2^* \circ \overline{d} \to j_3^*
$$
is equal to the map
$$
j_1^* \circ \overline{e} \to j_3^*
$$
via the identification $\overline{e} = \overline{c} \circ \overline{d}$.
This was proven in Lemma \href{https://stacks.math.columbia.edu/tag/0ATQ}{lemma compose base change maps (Tag 0ATQ)}
(note that in the current case the morphisms $f', g'$ of that
lemma are equal to $\text{id}_X$).
\end{proof}
\end{document}
